% Class and style of the document.
\documentclass[12pt]{article}

% Package and configuration.

% Language and encoding packages.
\usepackage[french]{babel}
\usepackage[utf8]{inputenc}
\usepackage[T1]{fontenc}

% Mathematics package.
\usepackage{amsmath}
\usepackage{amstext}

% Figures package.
\usepackage{graphicx}
\usepackage{wrapfig}

% Charts package.
\usepackage{pgfplots}
\pgfplotsset{compat=1.16}
\usepgfplotslibrary{units}
% Define style for charts.
\pgfplotsset{custom plot/.style={
        width            = 13cm,
        height           = 9cm,
        axis lines       = left,
        grid             = major,
        grid style       = {dashed},
        legend pos       = north west,
        ylabel           = Temps d'exécution,
        y unit           = \si{\second},
        enlarge y limits = upper,
        xlabel           = Taille des données,
        x unit           = N \text{ blocs}*K \text{ réels double précision},
        xticklabel style = {
            /pgf/number format/precision = 3,
        },
        log ticks with fixed point,
    }
}
\pgfplotsset{custom lines/.style={
        smooth,
        tension = 0.1,
    }
}
\pgfplotsset{custom line/.style={
        custom lines,
        color   = red,
        mark options = {
            fill = black,
        }
    }
}

% Symbols and numbers.
\usepackage{siunitx}
\usepackage{numprint}

% Formatter package.
%% Page break between section.
\usepackage{titlesec}
\newcommand{\sectionbreak}{\clearpage}
%% Use clickable URLs.
\usepackage{url}

% Hypertext links.
\usepackage[hidelinks]{hyperref}

% Footer.
\usepackage{lastpage}
\usepackage{fancyhdr}
\pagestyle{fancy}
% Clear headers and footers.
\fancyhead{}
% \renewcommand{\headrulewidth}{0pt}
% \renewcommand{\footrulewidth}{0pt}
% Define footers.
\fancyfoot[RE,LO]{Pierre AYOUB}
\fancyfoot[RE,CO]{\text{\thepage} sur \pageref{LastPage}}
\fancyfoot[RE,RO]{Océane FLAMANT}

% Paragraph.
\setlength{\parskip}{1ex}

% Quote.
\usepackage[style=french,french=guillemets]{csquotes}

\begin{document}

\title{Tri parallèle d'un tableau avec OpenMP}
\author{Pierre AYOUB -- Océane FLAMANT}

\maketitle

\begin{figure}[b]
    \centering
    \includegraphics[scale=0.3]{pictures/isty.jpg}
\end{figure}

\tableofcontents

\section{Introduction}

Durant la réalisation de ce projet, nous avions comme objectifs, dans un premier
temps, d'implémenter en langage C un algorithme de tri parallèle destiné à trier
une grande base de données. Dans un second temps, il nous fallait effectuer des
mesures du temps d'exécution de notre programme en fonction de plusieurs
paramètres, tels que le nombre de threads utilisés ou encore la taille des données.

Plusieurs contraintes nous étaient imposées :
\begin{itemize}  
    \item{Le squelette de l'algorithme à suivre était donné dans l'énoncé ;}
    \item{Nous avions à utiliser la bibliothèque de programmation parallèle OpenMP pour paralléliser notre algorithme;}
    \item{Nous avions trois fonctions principales à implémenter de la manière suivante :}
        \begin{description}
            \item[generator() :] {remplit un tableau de taille $K$ avec des nombres réels aléatoires;}
            \item[tri() :] {trie un tableau de taille $K$ ;}
            \item[tri-merge() :] {à partir de deux tableaux chacun de taille $K$ et triés, renvoie deux tableaux triés vérifiant que toutes les valeurs du premier tableau sont inférieures à celles du deuxième.}
        \end{description}
\end{itemize}

\section{Développement du programme}

Dans cette partie, nous allons expliquer et justifier les choix que nous avons
fait afin d'implémenter l'algorithme demandé, ainsi que les modifications qui
ont dû être effectuées.

\subsection{Structure du programme}

Nos fonctions \emph{generator()} et \emph{tri-merge()} suivent scrupuleusement
l'algorithme défini dans l'énoncé. Cependant, nous avions le libre arbitre
concernant l'implémentation de la fonction \emph{tri()}.
    
Pour l'algorithme de tri, nous avons choisi le \emph{tri à bulles} (connu aussi
sous le nom de \emph{tri par permutation} ou encore \emph{tri par propagation}).
Nous reviendrons sur ces performances plus tard, mais nous pouvons souligner le
fait que c'est un tri simple à implémenter et qu'il fait partie de la famille
des tris en place : il n'a pas besoin de recopier un tableau annexe pour que le
tri s'effectue. 

\subsection{Adaptation de l'algorithme}

Nous avons rencontré quelques déconvenues avec l'algorithme donné. En effet, une
fois nos fonctions réalisées et vérifiées, nous avons lancé le programme et le
résultat obtenu n'était pas le bon. Les valeurs contenues dans les tableaux
étaient triés mais les tableaux n'étais pas dans l'odre attendu. Nous avons donc
dû effectuer de légères adaptations.

Tout d'abord, en C, les indices commencent à $0$ et non à $1$, comme donné dans
l'algorithme donc pour chaque \emph{boucle for} commence à $0$ et s'arrête à 
$N-1$. Ensuite nous avons supprimé les $+1$ des variables $k$, $b1$ et $b2$ car 
les boucles commencent à 0. Enfin, pour $b2$, le $K$ est devenu $k$ (ce dernier
étant une correction d'une erreur de l'énoncé). 

En résumé :
\begin{itemize}
    \item $k = 1+(j\bmod 2) \Rightarrow k = (j\bmod 2)$
    \item $b_1 = 1+(k+2*i)\bmod N \Rightarrow b_1 = (k+2*i)\bmod N$
    \item $b_2 = 1+(K+2*i+1)\bmod N \Rightarrow b_2 = (k+2*i+1)\bmod N$
\end{itemize}

\section{Performances}

Nous allons maintenant discuter des performances de notre programme. D'un
point de vue algorithmique, dû à notre choix de \emph{tri à bulle}, les
performances ne devraient pas être excellentes. En effet, le \emph{tri à bulle}
a une complexité en moyenne en $\mathcal{O}(n^2)$, ce qui est bien moins bon
qu'un algorithme tel que le \emph{tri rapide} ayant une complexité en moyenne en
$\mathcal{O}(n.\log n)$. Cependant, étant donné que nous sommes dans une projet
de programmation parallèle et non d'algorithmique, il y a deux raisons qui
justifient le choix du \emph{tri à bulle} :
\begin{itemize}
    \item Une \enquote{mauvaise} complexité algorithmique nous permettra de mieux
        mettre en évidence l'importance du parallélisme ;
    \item L'algorithme parcours de multiples fois le tableau sur toute sa
        longueur, ainsi cela va rendre possible la mise en avant des défauts de
        cache.
\end{itemize}

\subsection{Mesures}

Avant de commenter nos mesures, il nous semble bon d'expliquer un minimum ce qui
va suivre. Notre programme voit son comportement modifié par 4 paramètres :
\begin{description}
    \item[Taille de la base de données :]{correspond au nombre de tableaux
        multiplié par la taille d'un tableau $(N*K)$. Il faut multiplier ce
        nombre par la taille d'un réel double précision pour l'exprimer en
        octet ;}
    \item[N :]{nombre de tableaux ;}
    \item[K :]{taille d'un tableau ;}
    \item[OMP\_NUM\_THREADS :]{nombre de threads exécutant le programme.}
\end{description}

Il est important de se rendre compte que si $N < \text{OMP\_NUM\_THREADS}$,
alors les $\text{OMP\_NUM\_THREADS} - N$ threads ne pourront pas obtenir de
travail.

De tous les paramètres des trois graphiques suivants, rien n'as été choisi au
hasard. Par défaut :
\begin{itemize}
    \item OMP\_NUM\_THREADS est égal aux nombre de c\oe{}urs de la machine hôte ;
        \item $N$ est calculé à partir d'une formule déterminée empiriquement,
            permettant d'approcher la valeur optimale de $N$ pour un $N*K$ donné
            après une série de mesures sur une de nos machines ;
        \item $K$ prend sa valeur en fonction de $N$.
\end{itemize}

Enfin, vous noterez que les échelles de nos graphiques sont des échelles
logarithmiques, avec un facteur 10 entre chaque graduation. 

\begin{figure}
    \begin{center}
        \begin{tikzpicture}
            \begin{loglogaxis} [
                    custom plot,
                    title = $N \text{ = } \exp(\log_{10}(N*K))$ ; $K \text{ = } \frac{(N*K)}{N}$ ; $\text{OMP\_NUM\_THREADS = } \mathopen|\text{CPU cores}\mathclose|$,
                ]
                \addplot+[custom line] table [
                    x       = data size,
                    y       = exec time,
                    col sep = comma,
                ] {charts/chart1.csv};
            \end{loglogaxis}
        \end{tikzpicture}
        \caption{Temps d'exécution en fonction de la taille des données}
        \label{time-to-data-size}
    \end{center}
\end{figure}

Comme nous pouvons le constater sur la Figure \ref{time-to-data-size}, le temps
d'exécution varie énormément en fonction de la taille des données, malgré que
les 3 paramètres présentés ci-dessus aient été choisis de manière optimale. En
dessous d'une taille de \numprint{100000} nombres, le temps d'exécution est
négligeable pour un utilisateur. Cependant, à partir d'une taille d'environ
\numprint{30000} nombres, nous dépassons définitivement le seuil des
\numprint{0.01}\si{\second} : cela s'explique par le fait que notre base de
données ne tient plus dans nos mémoires caches $L_1 \text{ et } L_2$. De plus, à
partir de ce seuil là se créé une droite de fonction linéaire, avec un
coefficient directeur environ égal à $1.5$. Cela paraît peu, mais cela n'est pas
à négliger puisque nous sommes sur des échelles logarithmiques : c'est-à-dire
que lorsque l'on multiplie la taille de nos données par $10$, le temps
d'exécution est environ multiplié par $30$. Le temps d'exécution va donc croître
très rapidement. Par exemple, pour \numprint{1000000} de nombres, nous étions à
$3\si{\second}$ d'exécution. A \numprint{10000000}, nous étions maintenant à
$100\si{\second}$. Nous pourrions largement imaginer qu'à \numprint{100000000},
nous serions à $3000\si{\second}$.

\begin{figure}
    \begin{center}
        \begin{tikzpicture}
            \begin{loglogaxis} [
                    custom plot,
                    title = $\text{OMP\_NUM\_THREADS = } \mathopen|\text{CPU cores}\mathclose|$ ; $K \text{ = } \frac{(N*K)}{N}$,
                ]
                \foreach \i in {2,16,64,256,1024}
                    \addplot+[custom lines] table [
                        x       = data size,
                        y       = exec time n eq \i,
                        col sep = comma,
                    ] {charts/chart2.csv};
                \addlegendentry{$N = 2$};
                \addlegendentry{$N = 16$};
                \addlegendentry{$N = 64$};
                \addlegendentry{$N = 256$};
                \addlegendentry{$N = 1024$};
            \end{loglogaxis}
        \end{tikzpicture}
        \caption{Temps d'exécution en fonction de la taille des données et du nombre de blocs}
        \label{time-to-data-size-and-blocks}
    \end{center}
\end{figure}

La Figure \ref{time-to-data-size-and-blocks} nous permet de mesurer l'influence
du paramètre $N$ sur le temps d'exécution de notre programme. Notez tout d'abord
que si $N > N*K$, alors le test ne peut pas avoir lieu : nous avons donc des
courbes qui ne commencent pas à une taille de $10$. Nous observons que pour des
petites tailles de base de données $(< 100)$, un $N$ petit est optimal.
Cependant, un $N$ petit montre très vite ses limites et est dépassé par des
valeurs de $N$ supérieures. En revanche, les grands $N$, dû aux changements
constant de tableau en cours de traitement, sont nettement moins adaptés pour
des petites tailles de base de données (le temps d'exécution peut être multiplé
par $10$ par rapport à un $N$ petit), mais devient très vite optimal pour de
grandes bases de données $(> \numprint{100000})$. Les différences les plus
extrêmes de notre graphique se situent à une taille de base de données de
\numprint{1000000}, où un $N=2$ donne un temps d'exécution de
\numprint{1000}\si{\second}, tandis qu'un $N=1024$ donne un temps d'exécution de
\numprint{3}\si{\second} : c'est une différence d'un facteur d'environ $333$ !

\begin{figure}
    \begin{center}
        \begin{tikzpicture}
            \begin{loglogaxis} [
                    custom plot,
                    legend pos = north west,
                    title = $N \text{ = } \exp(\log_{10}(N*K))$ ; $K \text{ = } \frac{(N*K)}{N}$,
                ]
                \foreach \i in {1,4,16,64}
                    \addplot+[custom lines] table [
                        x       = data size,
                        y       = exec time omp_num_threads eq \i,
                        col sep = comma,
                    ] {charts/chart3.csv};
                \addlegendentry{$\text{OMP\_NUM\_THREADS} = 1$};
                \addlegendentry{$\text{OMP\_NUM\_THREADS} = 4$};
                \addlegendentry{$\text{OMP\_NUM\_THREADS} = 16$};
                \addlegendentry{$\text{OMP\_NUM\_THREADS} = 64$};
            \end{loglogaxis}
        \end{tikzpicture}
        \caption{Temps d'exécution en fonction de la taille des données et du nombre de threads}
        \label{time-to-data-size-and-threads}
    \end{center}
\end{figure}

La Figure \ref{time-to-data-size-and-threads} nous informe de la différence du
temps d'exécution en fonction du nombre de threads se chargeant du tri. Comme
attendu, un grand nombre de threads sur un petit nombre de données augmente le
temps d'exécution dû aux nombreux changements de contextes inutiles. Cependant,
avoir plusieurs threads devient bénéfique à partir d'une taille comprise entre
\numprint{1000} et \numprint{10000} nombres. En revanche, nous pouvons remarquer
que le temps d'exécution est similaire entre $4$ et $64$ threads à partir d'une
taille d'\numprint{1000000} de données : la cause étant que la machine ne
disposant que de $4$ c\oe{}urs. Si nous avions un CPU disposant de plus de
c\oe{}urs, alors nous aurions vu le temps d'exécution continuer de diminuer en
augmentant le nombre de threads. La différence entre $1$ et $4$ threads pour
\numprint{1000000} de nombres est tout de même non-négligeable, passant
respectivement d'un temps d'exécution de $200\si{\second}$ à $100\si{\second}$.
Nous pourrions nous attendre à un facteur $4$, mais nous n'observons qu'un
facteur $2$ : l'origine de ce phénomène vient de \emph{l'HyperThreading},
exposant au système d'exploitation $4$ c\oe{}urs logiques tandis que le CPU ne
possède seulement que $2$ c\oe{}urs physiques.

\subsection{Perspectives d'amélioration}
    
Plusieurs moyens peuvent être mis en \oe{}uvre pour améliorer les performances
de notre programme :
\begin{itemize}
    \item Augmenter le degré de parallélisme : paralléliser de nouvelles
        boucles, découper notre tableau sur de nouveaux niveaux.
    \item Améliorer la gestion de la mémoire cache : effectuer des
        benchmarks pour déterminer la taille précise à partir de laquelle
        nos \emph{cache miss} augmentent de manière notable.
    \item Améliorer la complexité algorithmique : utiliser un algorihtme de
        tri plus efficace (tel que le \emph{tri fusion}, le \emph{tri
        rapide} ou encore le \emph{tri par tas}).
\end{itemize}

\section{Conclusion}

Les mesures de performance de notre programme auront été une tâche intéréssante,
car plusieurs problématiques ont été introduites pour y répondre convenablement.
Certains choix n'auront pas été de toute simplicité, tel que le choix de l'échelle
des axes, du nombre de graphiques, des échantillons représentatifs, etc. De plus,
nous avons créé un script shell permettant de recréer et/ou modifier nos
expériences de manière très simple.

Concernant la partie analyse, il a été instructif de comprendre les correlations
entre les différences de performance induites par chaque paramètre. Nous pouvons
aussi noter à quel point, sur une machine donnée, un paramètre peut-être bien
plus important qu'un autre (ici, $N$ est bien plus influant que
OMP\_NUM\_THREADS).

Nous pouvons donc en conclure que pour obtenir un programme performant, il faut
à la fois une étude algorithmique poussée pour obtenir la meilleure complexité,
mais aussi prendre en compte les caractéristiques de l'algorithme qui pourraient
influencer le matériel.

\end{document}
